% ============================================
% 术语格式参考 (Glossary Format Reference)
% 中英对照格式：中文 (English)
% ============================================

% --------------------------------------------
% 行内术语格式
% --------------------------------------------

% 首次出现时使用完整格式
人工智能 (Artificial Intelligence, AI) 是计算机科学的一个重要分支。

% 后续出现可只用缩写
AI 技术在近年来取得了显著进展。


% --------------------------------------------
% 常见学术术语示例
% --------------------------------------------

% 研究方法类
定性研究 (Qualitative Research)
定量研究 (Quantitative Research)
实证分析 (Empirical Analysis)
案例研究 (Case Study)

% 统计分析类
显著性水平 (Significance Level)
置信区间 (Confidence Interval)
回归分析 (Regression Analysis)
相关系数 (Correlation Coefficient)

% 通用学术类
文献综述 (Literature Review)
理论框架 (Theoretical Framework)
研究假设 (Research Hypothesis)
自变量 (Independent Variable)
因变量 (Dependent Variable)

% C 语言基础类
指针 (Pointer)
数组 (Array)
结构体 (Structure)
联合体 (Union)
枚举 (Enumeration)
宏定义 (Macro Definition)
预处理器 (Preprocessor)
头文件 (Header File)
源文件 (Source File)
编译器 (Compiler)
链接器 (Linker)

% C 语言内存管理类
内存分配 (Memory Allocation)
动态内存 (Dynamic Memory)
栈 (Stack)
堆 (Heap)
内存泄漏 (Memory Leak)
野指针 (Wild Pointer)
悬空指针 (Dangling Pointer)
空指针 (Null Pointer)
内存对齐 (Memory Alignment)
缓冲区溢出 (Buffer Overflow)

% C 语言指针操作类
解引用 (Dereference)
取地址 (Address-of)
指针运算 (Pointer Arithmetic)
函数指针 (Function Pointer)
指针数组 (Array of Pointers)
数组指针 (Pointer to Array)
多级指针 (Multi-level Pointer)
空悬指针 (Dangling Pointer)
常量指针 (Constant Pointer)
指针常量 (Pointer to Constant)

% C 语言数据结构类
链表 (Linked List)
单向链表 (Singly Linked List)
双向链表 (Doubly Linked List)
循环链表 (Circular Linked List)
二叉树 (Binary Tree)
哈希表 (Hash Table)
队列 (Queue)
栈 (Stack)


% --------------------------------------------
% 术语表格式（用于附录）
% --------------------------------------------
\begin{table}[H]
  \centering
  \caption{主要术语对照表}
  \label{tab:glossary}
  \begin{tabular}{lll}
    \toprule
    \textbf{中文} & \textbf{英文} & \textbf{缩写} \\
    \midrule
    人工智能 & Artificial Intelligence & AI \\
    机器学习 & Machine Learning & ML \\
    深度学习 & Deep Learning & DL \\
    自然语言处理 & Natural Language Processing & NLP \\
    \bottomrule
  \end{tabular}
\end{table}
